\chapter{\ifenglish Introduction\else บทนำ\fi}

\section{\ifenglish Project rationale\else ที่มาของโครงงาน\fi}

ในปัจจุบันอุตสาหกรรมเกมดิจิทัลมีการเติบโตอย่างรวดเร็ว และมีบทบาทสำคัญทั้งในด้านความบันเทิง การศึกษา และเทคโนโลยีสารสนเทศ การพัฒนาเกมจึงเป็นการประยุกต์ใช้ความรู้จากหลายสาขา เช่น การเขียนโปรแกรมคอมพิวเตอร์ การออกแบบระบบซอฟต์แวร์ การออกแบบกราฟิก และการออกแบบประสบการณ์ผู้ใช้

โครงงานนี้มีวัตถุประสงค์เพื่อพัฒนาเกมต้นแบบชื่อว่า \textbf{Marionette} ซึ่งเป็นเกมแนว 2 มิติประเภท Action RPG ที่ผสมผสานรูปแบบการเล่นแบบ Casual และ Mini-game ภายในธีมของโรงละครหุ่นเชิด ผู้เล่นจะรับบทเป็นตัวละครที่ถูกจับมาแปลงเป็นหุ่นเชิดและต้องพยายามหาทางหลบหนีเพื่อกลับคืนสู่อิสรภาพ

เกมถูกออกแบบให้มีรูปแบบการเล่นที่หลากหลาย เช่น การสำรวจแผนที่เพื่อเก็บไอเทม การต่อสู้กับศัตรู การแสดงบนเวทีในรูปแบบมินิเกม และการต่อสู้กับบอสในแต่ละด่าน เพื่อสร้างประสบการณ์การเล่นที่มีความท้าทายและหลากหลาย โครงงานนี้จึงเป็นการศึกษาและพัฒนากระบวนการสร้างเกมตั้งแต่ขั้นตอนการออกแบบระบบเกมไปจนถึงการพัฒนาต้นแบบเกมที่สามารถเล่นได้จริง


\section{\ifenglish Objectives\else วัตถุประสงค์ของโครงงาน\fi}
\begin{enumerate}
    \item เพื่อออกแบบและพัฒนาเกมต้นแบบแนว 2 มิติประเภท Action RPG
    \item เพื่อศึกษากระบวนการพัฒนาเกมตั้งแต่การออกแบบระบบเกมจนถึงการพัฒนาโปรแกรม
    \item เพื่อพัฒนาระบบหลักของเกม เช่น ระบบการต่อสู้ ระบบช่องเก็บของ และระบบศัตรู
    \item เพื่อสร้างเกมต้นแบบที่สามารถเล่นได้จริงในระดับด่านแรก
    \item เข้าร่วมแข่งขันในงาน NSC 2026 (National Software Contest) และได้รับการยอมรับในระดับประเทศ
\end{enumerate}

\section{\ifenglish Project scope\else ขอบเขตของโครงงาน\fi}

\subsection{\ifenglish Hardware scope\else ขอบเขตด้านฮาร์ดแวร์\fi}

โครงงานนี้พัฒนาเกมสำหรับการทำงานบนคอมพิวเตอร์ส่วนบุคคล โดยใช้ฮาร์ดแวร์มาตรฐานที่สามารถรองรับการพัฒนาและทดสอบเกมได้ เช่น

\begin{itemize}
    \item คอมพิวเตอร์ส่วนบุคคล
    \item หน้าจอแสดงผล 
    \item เมาส์และคีย์บอร์ดสำหรับควบคุมการเล่นเกม
\end{itemize}

\subsection{\ifenglish Software scope\else ขอบเขตด้านซอฟต์แวร์\fi}

ซอฟต์แวร์ในโครงงานนี้เป็นเกมแนว 2 มิติที่พัฒนาโดยใช้ Unity Game Engine โดยมีระบบหลักของเกม ได้แก่
\begin{itemize}
    \item ระบบการสำรวจแผนที่
    \item ระบบการต่อสู้กับศัตรู
    \item ระบบไอเทมและอุปกรณ์
    \item ระบบมินิเกมสำหรับการแสดงบนเวที
    \item ระบบบอสของแต่ละด่าน
\end{itemize}

\section{\ifenglish Expected outcomes\else ประโยชน์ที่ได้รับ\fi}

\begin{itemize}
    \item ได้ต้นแบบเกมที่สามารถเล่นได้จริง
    \item ได้ความรู้เกี่ยวกับกระบวนการพัฒนาเกม
    \item ได้พัฒนาทักษะด้านการเขียนโปรแกรมและการออกแบบระบบ
    \item สามารถนำความรู้ไปประยุกต์ใช้ในการพัฒนาเกมหรือซอฟต์แวร์อื่นในอนาคต
\end{itemize}

\section{\ifenglish Technology and tools\else เทคโนโลยีและเครื่องมือที่ใช้\fi}

\subsection{\ifenglish Hardware technology\else เทคโนโลยีด้านฮาร์ดแวร์\fi}

\begin{itemize}
    \item Personal Computer
    \item Keyboard
    \item Mouse
\end{itemize}

\subsection{\ifenglish Software technology\else เทคโนโลยีด้านซอฟต์แวร์\fi}

\begin{itemize}
    \item Unity Game Engine
    \item Visual Studio / Visual Studio Code / Antigravity
    \item C\# Programming Language
    \item Generative AI Tools (e.g., ChatGPT, Gemini) for content creation and development assistance
\end{itemize}

\section{\ifenglish Project plan\else แผนการดำเนินงาน\fi}

\begin{plan}{11}{2025}{5}{2026}
    \planitem{11}{2025}{12}{2025}{ศึกษาข้อมูลและออกแบบ Game Design}
    \planitem{1}{2026}{1}{2026}{นำเสนอต่อกรรมการและปรับปรุงตามคำแนะนำ}
    \planitem{2}{2026}{2}{2026}{พัฒนาเกมต้นแบบด่านแรก}
    \planitem{3}{2026}{3 }{2026}{Prototype ด่านแรก}
    \planitem{4}{2026}{4 }{2026}{แก้ไขและปรับปรุงเกมต้นแบบตามผลการทดสอบ}
    \planitem{5}{2026}{5}{2026}{ยื่นข้อเสนอโครงการส่ง NSC 2026}
\end{plan}


\section{\ifenglish Roles and responsibilities\else บทบาทและความรับผิดชอบ\fi}

ในการดำเนินโครงงานนี้ สมาชิกในทีมได้มีการแบ่งหน้าที่ความรับผิดชอบตามความถนัดและความเชี่ยวชาญของแต่ละบุคคล เพื่อให้การพัฒนาโครงงานสามารถดำเนินไปได้อย่างมีประสิทธิภาพ โดยมีการแบ่งบทบาทหน้าที่ดังแสดงในตารางด้านล่าง

\begin{center}
\begin{tabular}{|c|c|c|p{6cm}|}
\hline
ลำดับ & ชื่อ - นามสกุล & รหัสนักศึกษา & หน้าที่ความรับผิดชอบ \\
\hline
1 & นายกานต์ ปราศัย & 660610742 & พัฒนาระบบการต่อสู้ภายในเกม ออกแบบและพัฒนาระบบบอส รวมถึงการจัดการแอนิเมชันของตัวละคร เช่น การทำ Rig และการควบคุมการเคลื่อนไหวของตัวละคร \\
\hline
2 & นายณัฐดนัย กำธรกิตติกุล & 660610755 & พัฒนาระบบการสำรวจแผนที่ภายในเกม รวมถึงการออกแบบและพัฒนาระบบไอเทมที่ผู้เล่นสามารถเก็บและใช้งานได้ระหว่างการเล่นเกม \\
\hline
3 & นายปัณณทัต เงินงาม & 660610775 & พัฒนาและออกแบบระบบมินิเกมที่ใช้ในช่วงการแสดงบนเวที ซึ่งเป็นหนึ่งในกลไกสำคัญของรูปแบบการเล่นภายในเกม \\
\hline
\end{tabular}
\end{center}

นอกจากนี้ สมาชิกทุกคนยังมีส่วนร่วมในการออกแบบแนวคิดของเกม การทดสอบระบบ (Testing) และการปรับปรุงแก้ไขข้อผิดพลาด (Debugging) เพื่อให้เกมต้นแบบสามารถทำงานได้อย่างสมบูรณ์

\section{งบประมาณ}
โครงงานนี้มีค่าใช้จ่ายที่เกี่ยวข้องกับการพัฒนาเกมต้นแบบ โดยส่วนใหญ่เป็นค่าใช้จ่ายด้านทรัพยากรดิจิทัลและบริการที่ใช้ในการพัฒนาเกม อาทิ ค่าจ้างนักวาดภาพสำหรับออกแบบตัวละครและองค์ประกอบกราฟิกของเกม ค่าใช้จ่ายในการจัดหาเพลงประกอบเกม รวมถึงค่าใช้จ่ายในการจัดซื้อทรัพยากรจาก Unity Asset Store ที่นำมาใช้สนับสนุนการพัฒนาระบบต่าง ๆ ภายในเกม นอกจากนี้ยังมีค่าใช้จ่ายสำหรับเครื่องมือที่ช่วยสนับสนุนกระบวนการพัฒนา เช่น บริการปัญญาประดิษฐ์ ที่ใช้ช่วยในการเขียนและปรับปรุงโค้ด เพื่อเพิ่มประสิทธิภาพและความรวดเร็วในการพัฒนาโครงงาน

\section{\ifenglish%
Impacts of this project on society, health, safety, legal, and cultural issues
\else%
ผลกระทบด้านสังคม สุขภาพ ความปลอดภัย กฎหมาย และวัฒนธรรม
\fi}

โครงงานนี้เป็นการพัฒนาเกมเพื่อการศึกษาและความบันเทิง ซึ่งไม่มีผลกระทบด้านความปลอดภัยหรือกฎหมายโดยตรง อย่างไรก็ตาม การพัฒนาเกมต้องคำนึงถึงเนื้อหาที่เหมาะสมต่อผู้เล่น รวมถึงการออกแบบเกมที่ไม่ส่งเสริมพฤติกรรมที่ไม่เหมาะสม นอกจากนี้ ความรู้ที่ได้จากโครงงานสามารถนำไปประยุกต์ใช้หรือสร้างแรงบันดาลใจในการพัฒนาซอฟต์แวร์หรือเกมอื่น ๆ ที่มีประโยชน์ต่อสังคมในอนาคต